\documentclass[a4paper,9pt]{extarticle}

%=========================================================
% PACKAGES
%=========================================================

\usepackage[
a4paper,
top=1.8cm,
bottom=1.5cm,
footskip=7mm,
left=1cm,
right=1cm,
headheight=14pt,
headsep=5mm
]{geometry}
\usepackage[T1]{fontenc}
\usepackage[utf8]{inputenc}
\usepackage[french]{babel}

\usepackage[sfdefault]{FiraSans}
\renewcommand{\familydefault}{\sfdefault}

\usepackage{multicol}
\setlength{\columnsep}{0.7cm}

\usepackage{amsmath,amssymb}

\usepackage{titlesec}
\usepackage{enumitem}

\usepackage{tcolorbox}
\tcbuselibrary{skins,breakable}

\usepackage{xcolor}

\usepackage{graphicx}
\usepackage{float}

\usepackage{tikz}

\usepackage{fancyhdr}
\usepackage[hidelinks]{hyperref}


%=========================================================
% PAGE SETTINGS
%=========================================================

\setlength{\parindent}{0pt}
\setlength{\parskip}{1.5pt}

\setlength{\abovedisplayskip}{4pt}
\setlength{\belowdisplayskip}{4pt}

\setlist[itemize]{
leftmargin=*,
itemsep=1pt,
topsep=1pt
}


%=========================================================
% COLORS
%=========================================================

\definecolor{Blue}{HTML}{E3F0FF}
\definecolor{Green}{HTML}{E3F7EA}
\definecolor{Pink}{HTML}{FFE5ED}
\definecolor{Orange}{HTML}{FFF0DC}
\definecolor{Purple}{HTML}{EEE7FF}
\definecolor{Yellow}{HTML}{FFFADB}
\definecolor{Teal}{HTML}{E0F5F3}


\definecolor{BlueTitle}{HTML}{285B8F}
\definecolor{GreenTitle}{HTML}{287052}
\definecolor{PinkTitle}{HTML}{A84063}
\definecolor{OrangeTitle}{HTML}{A96500}
\definecolor{PurpleTitle}{HTML}{6245A5}
\definecolor{TealTitle}{HTML}{1F6F68}


%=========================================================
% HEADER
%=========================================================

\pagestyle{fancy}

\fancyhf{}

\fancyhead[L]{\small\textbf{Modèles stochastiques}}
\fancyhead[C]{\small Introduction et variables aléatoires}
\fancyhead[R]{\small Week 01 --- 07-13/09/2026}

\fancyfoot[C]{\small\href{https://github.com/AlphaXT38G/notes_ba5}{github.com/AlphaXT38G/notes\_ba5}}
\fancyfoot[R]{\small \thepage}

\renewcommand{\headrulewidth}{0.4pt}



%=========================================================
% TITLES
%=========================================================

\titleformat{\section}
{\large\bfseries\color{BlueTitle}}
{}{0pt}{}

\titlespacing{\section}
{0pt}{6pt}{3pt}


% separator

\newcommand{\separator}{
\vspace{2pt}
\hrule
\vspace{3pt}
}


%=========================================================
% BOXES
%=========================================================

% Boxes are kept whole (unbreakable): inside multicols, a box split by tcolorbox
% leaves an empty column behind it.

% DEFINITION - blue
\newtcolorbox{definition}{
enhanced,
unbreakable,
colback=Blue,
colframe=BlueTitle,
boxrule=0.7pt,
arc=0mm,
title=\textbf{Definition},
fonttitle=\bfseries
}

% IMPORTANT - pink
\newtcolorbox{important}{
enhanced,
unbreakable,
colback=Pink,
colframe=PinkTitle,
boxrule=0.7pt,
arc=0mm,
title=\textbf{Key idea},
fonttitle=\bfseries
}

% FORMULAS - green
\newtcolorbox{formula}{
enhanced,
unbreakable,
colback=Green,
colframe=GreenTitle,
boxrule=0.7pt,
arc=0mm,
title=\textbf{Key formulas},
fonttitle=\bfseries
}

% WARNING - orange
\newtcolorbox{warning}{
enhanced,
unbreakable,
colback=Orange,
colframe=OrangeTitle,
boxrule=0.7pt,
arc=0mm,
title=\textbf{Warning},
fonttitle=\bfseries
}

% EXAMPLE - purple
\newtcolorbox{example}{
enhanced,
unbreakable,
colback=Purple,
colframe=PurpleTitle,
boxrule=0.7pt,
arc=0mm,
title=\textbf{Example},
fonttitle=\bfseries
}

% SUMMARY - orange
\newtcolorbox{summary}{
enhanced,
unbreakable,
colback=Yellow,
colframe=OrangeTitle,
boxrule=0.7pt,
arc=0mm,
title=\textbf{Summary},
fonttitle=\bfseries
}



% SIMPLE - teal: the same idea again, in the plainest words, with an everyday picture
\newtcolorbox{simple}{
enhanced,
unbreakable,
colback=Teal,
colframe=TealTitle,
boxrule=0.7pt,
arc=0mm,
title=\textbf{Expliqué simplement},
fonttitle=\bfseries
}



%=========================================================
% FIGURE BOX
%=========================================================

\newtcolorbox{figurebox}{
enhanced,
colback=white,
colframe=gray,
boxrule=0.5pt,
arc=0mm,
center
}



%=========================================================
% DOCUMENT
%=========================================================

\begin{document}

%\footnotesize   --> To reduce the size of the text

\thispagestyle{fancy}



\begin{multicols}{2}


%=========================================================
\section{Organisation du cours}


\begin{summary}

\begin{itemize}

\item Enseignant : Patrick Thiran. Assistants : Anthony Bardou, Raphaël Fua.

\item Cours en présentiel uniquement : mardi 10h15 (CM3) et vendredi 13h15 (CM2). Les
enregistrements (mediaspace) sont un \textbf{complément}, ils ne remplacent pas le cours.

\item Exercices le mardi et le vendredi, dans la même salle que le cours qui précède.

\item \textbf{Midterm} (modules 1 à 3) : vendredi 30 octobre 2026, 13h15, \textbf{30 pts}.
Une partie QCM (7 questions : $+2$ si juste, $-1$ si faux, $0$ sans réponse) et une partie
questions ouvertes.

\item \textbf{Examen final} en session (modules 1 à 8) : \textbf{70 pts}.

\item Seul document autorisé : le \textbf{formulaire officiel} imprimé, sans annotations
(sauf surlignage).

\end{itemize}

\end{summary}


\textbf{Plan du semestre}

\begin{itemize}

\item Rappels de proba : variable aléatoire (semaine 1)

\item Rappels de proba : vecteurs et suites de v.a. (semaines 2--3)

\item Processus stochastiques à temps continu : 2nd ordre (semaines 4--6)

\item Processus stochastiques à temps discret : 2nd ordre (semaines 6--7)

\item Processus de Poisson (semaines 8--9)

\item Chaînes de Markov à temps discret puis continu, files d'attente (semaines 10--14)

\end{itemize}


\begin{important}

Questions ouvertes : réponse correcte et complète = tous les points, réponse fausse ou
inintelligible = 0, réponse correcte mais incomplète = crédit partiel possible.

\end{important}


\begin{simple}

Ce cours, c'est apprendre à décrire des choses qui bougent \textbf{au hasard} (du bruit sur un
fil, des clients qui arrivent, des paquets sur un réseau) avec des maths. Cette semaine, on
révise la brique de base : \textbf{une seule} variable aléatoire.

\end{simple}

\end{multicols}


\end{document}
